\documentclass[twocolumn]{aastex631}
\begin{document}
\title{The reionization ionizing-photon-budget shortfall for star-forming galaxies is not robust to systematics at z\~{}5 (o32 systematic corner)}
\author{NebulaMind Lab (autonomous pipeline)}
\affiliation{NebulaMind Lab, \url{https://lab.nebulamind.net}}
\begin{abstract}
Reionization ionizing-photon-budget reconciliation at z\~{}5: star-forming galaxies require f\_esc=0.025 (+0.025/-0.013) to close the budget (Madau-Dickinson SFRD, log xi\_ion=25.5+/-0.15, clumping C=2-5, JWST-SFRD tail), versus indirect-proxy-inferred f\_esc=0.080 (+0.146/-0.051) from LzLCS O32/beta calibrations. Median delta(required-inferred)=-0.052 dex-frac (16-84\%: -0.195 to +0.002); 17\% of the systematic MC shows a shortfall. Conclusion: the budget CLOSES within the systematic, and the sign holds under both O32 and beta calibrations. The result is bounded by the xi\_ion x clumping x proxy-calibration systematic, not by statistics. Generated autonomously from public data (jwst) via the NebulaMind Lab runner: a bounded, reproducible, descriptive study.
\end{abstract}
\section{Introduction}
Recent studies have highlighted a potential crisis in the reionization photon budget, suggesting that current estimates of star-forming galaxies' ionizing photon production may not be sufficient to account for the observed reionization process [Muñoz2024]. This discrepancy has sparked interest in reconciling the ionizing photon budget using literature-anchored calculations. Previous work has explored various aspects of this problem, including the role of galaxy ionizing photon budgets [Duncan2015] and the challenges posed by absorption-dominated reionization scenarios [Davies2021].
\section{Data and method}
To address this issue, our study employs a systematic reconciliation approach based on published literature values. We utilize the Madau \& Dickinson (2014) analytic fitting function for the cosmic star formation rate density (SFRD), along with adopted values for xi\_ion and O32/beta f\_esc proxy calibrations from LzLCS [Chisholm+22, Flury+22; Simmonds+24]. Our method focuses on calculating the ionizing-photon-budget at z\~{}5 using these parameters.
\section{Result}
Our analysis reveals that star-forming galaxies require an escape fraction of f\_esc=0.025 (+0.025/-0.013) to reconcile the reionization ionizing-photon-budget, assuming a Madau-Dickinson SFRD, log xi\_ion=25.5+/-0.15, and clumping factor C=2-5. In contrast, indirect-proxy-inferred f\_esc values from LzLCS O32/beta calibrations yield f\_esc=0.080 (+0.146/-0.051). The median delta between the required and inferred escape fractions is -0.052 dex-frac (16-84\%: -0.195 to +0.002), with 17\% of systematic Monte Carlo simulations showing a shortfall.
\begin{figure}\centering\includegraphics[width=\columnwidth]{result.png}\caption{The reionization ionizing-photon-budget shortfall for star-forming galaxies is not robust to systematics at z\~{}5 (o32 systematic corner)}\end{figure}
\section{Caveats}
It is essential to acknowledge the limitations of our approach, which relies on automated, single-selection, uncalibrated measurements. The accuracy of our results depends heavily on the adopted literature values and calibrations, which may introduce uncertainties due to variations in observational data and modeling assumptions. Additionally, our method does not account for potential systematic errors or biases inherent in the published studies we draw upon. Further research is needed to refine these estimates and address the complexities of reionization photon budget calculations.
\section*{References}
\footnotesize
[Muoz2024] Muñoz et al. (2024). Reionization after JWST: a photon budget crisis?. bibcode 2024MNRAS.535L..37M \par [Park2022] Park et al. (2022). Calibrating excursion set reionization models to approximately conserve ionizing photons. bibcode 2022MNRAS.517..192P \par [Duncan2015] Duncan et al. (2015). Powering reionization: assessing the galaxy ionizing photon budget at z < 10. bibcode 2015MNRAS.451.2030D \par [Davies2021] Davies et al. (2021). The Predicament of Absorption-dominated Reionization: Increased Demands on Ionizing Sources. bibcode 2021ApJ...918L..35D \par [Madau2017] Madau (2017). Cosmic Reionization after Planck and before JWST: An Analytic Approach. bibcode 2017ApJ...851...50M

\end{document}
